\chapter{\texorpdfstring{\textit{Prompts}}{Prompts} Utilizados nos Experimentos}\label{ap:prompts}

Este apêndice reproduz integralmente os três modelos de \textit{prompt} empregados nas rodadas experimentais descritas no Capítulo~\ref{cap:metodologia}. As variáveis entre chaves (\texttt{\{language\}}, \texttt{\{problem\}}, \texttt{\{time\_limit\}} e \texttt{\{memory\_limit\}}) são preenchidas dinamicamente pelo módulo \texttt{PromptBuilder} no momento de cada submissão.

\section*{A.1 \textit{Prompt Zero-Shot} (\texttt{zero\_shot.txt})}

\begin{small}
\VerbatimInput[breaklines, breakanywhere, fontsize=\footnotesize, frame=single, framesep=2mm, rulecolor=\color{black!40}]{Apendices/prompts/zero_shot.txt}
\end{small}

\section*{A.2 \textit{Prompt Few-Shot} em C (\texttt{few\_shot\_c.txt})}

\begin{small}
\VerbatimInput[breaklines, breakanywhere, fontsize=\footnotesize, frame=single, framesep=2mm, rulecolor=\color{black!40}]{Apendices/prompts/few_shot_c.txt}
\end{small}

\section*{A.3 \textit{Prompt Chain-of-Thought} -- Metodologia de Pólya (\texttt{chain\_of\_thought\_polya.txt})}

\begin{small}
\VerbatimInput[breaklines, breakanywhere, fontsize=\footnotesize, frame=single, framesep=2mm, rulecolor=\color{black!40}]{Apendices/prompts/chain_of_thought_polya.txt}
\end{small}
